%
% polstellen.tex -- Lage der Polstellen in Abhängigkeit von \varepsilon
%
% (c) 2021 Prof Dr Andreas Müller, OST Ostschweizer Fachhochschule
%
\documentclass[tikz]{standalone}
\usepackage{amsmath}
\usepackage{times}
\usepackage{txfonts}
\usepackage{pgfplots}
\usepackage{csvsimple}
\usetikzlibrary{arrows,intersections,math,calc}
\definecolor{darkred}{rgb}{0.8,0,0}
\definecolor{darkgreen}{rgb}{0,0.6,0}
\begin{document}
\def\skala{1}
\begin{tikzpicture}[>=latex,thick,scale=\skala]

\draw[->] (-8.5,0) -- (3.2,0) coordinate[label={$\operatorname{Re}$}];
\draw[->] (0,-2.9) -- (0,2.9) coordinate[label={left:$\operatorname{Im}$}];


\draw[color=blue] (0,0) circle[radius=2];
\node[color=blue] at (1.52,1.68) {$C$};

%\draw (-2,-0.14) -- (-2,0.14);
\fill[color=blue] (-2,0) circle[radius=0.07];
\node[color=blue,anchor=south east] at (-1.92,0.1) {$-1$};


\fill[color=darkred] (-0.6667,0) circle[radius=0.07];
\node[color=darkred] at (-0.6667,0.32) {$z_+$};
\fill[color=darkred] (-6,0) circle[radius=0.07];
\node[color=darkred] at (-6,0.32) {$z_-$};


\draw[->, color=darkred, line width=1.1pt] (-6.3,0) -- (-8.2,0);
\node[color=darkred, anchor=north] at (-7.15,-0.00) {$\varepsilon \to 0^+$};
\draw[->, color=darkred, line width=1.1pt] (-5.7,0) -- (-2.3,0);
\node[color=darkred, anchor=north] at (-4.0,-0.00) {$\varepsilon \to 1^-$};


\draw[->, color=darkred, line width=1.1pt] (-0.9,0) -- (-1.95,0);
\node[color=darkred, anchor=north] at (-1.35,-0.00) {$\varepsilon \to 1^-$};
\draw[->, color=darkred, line width=1.1pt] (-0.45,0) -- (-0.05,0);
\node[color=darkred, anchor=north west] at (-0.35,-0.00) {$\varepsilon \to 0^+$};

\end{tikzpicture}
\end{document}

